\documentclass[11pt]{article}
\usepackage[margin=1in]{geometry}
\usepackage{amsmath,amssymb,booktabs,longtable,array}
\usepackage[T1]{fontenc}
\usepackage[utf8]{inputenc}
\title{Two-Group Correctness Attribution Audit}
\author{Codex Local Audit}
\date{}
\begin{document}
\maketitle

\section{Task Background and Problem Statement}
This report re-audits two already completed result lines:
\begin{itemize}
\item the closed reference group \texttt{10.4.1.31},
\item the controlled portability group \texttt{194.1.1.1}.
\end{itemize}
The point of the re-audit is not to rerun portability or extend the workflow, but to separate three logically distinct possibilities:
\[
  \text{(i) BS layer error}, \qquad
  \text{(ii) AI / completeness error}, \qquad
  \text{(iii) quotient interpretation error}.
\]
The warning sign forcing this split is the stage-2 SG 194 result
\[
  BS/AI \cong \mathbb{Z}^{16},
\]
which is too obviously free-dominated to be read naively as a standard finite indicator group.

\section{Current Results Under Review}
The audited headline outputs are:
\begin{align*}
  BS_{10,\mathrm{single}}/AI_{10,\mathrm{single}} &\cong \mathrm{Z}_{2} \times \mathrm{Z}_{2}, \\
  BS_{10,\mathrm{double}}/AI_{10,\mathrm{double}} &\cong \mathrm{Z}^{2} \times \mathrm{Z}_{2} \times \mathrm{Z}_{2} \times \mathrm{Z}_{2} \times \mathrm{Z}_{2}, \\
  BS_{194,\mathrm{single}}/AI_{194,\mathrm{single}} &\cong \mathrm{Z}^{16}, \\
  BS_{194,\mathrm{double}}/AI_{194,\mathrm{double}} &\cong \mathrm{Z}^{16}.
\end{align*}

\section{BS Formalism and Audit Method}
In all four cases the k-space layer is represented as an exact compatibility problem
\[
  C n = 0,
  \qquad
  BS = \ker_{\mathbb{Z}}(C).
\]
The audit checks:
\begin{enumerate}
\item whether geometry / manifolds / unknown ordering are explicit,
\item whether the matrix shape, rank, and nullity are consistent with the stored basis files,
\item whether plane necessity is supported by an added-rank test where applicable,
\item whether the kernel basis count matches the reported BS rank.
\end{enumerate}

\section{AI / Completeness Formalism and Audit Method}
The atomic lattice is treated as an integral span
\[
  AI = \operatorname{span}_{\mathbb{Z}}\{a_1,a_2,\dots\},
\]
generated by induced local objects. The audit does not simply trust ``complete'' flags. Instead it checks:
\begin{enumerate}
\item whether the family table is explicitly closed,
\item whether the local site-symmetry census is mathematically constrained enough to make missing local objects unlikely,
\item whether the stored AI basis coefficients reproduce the quoted quotient under a fresh Smith decomposition.
\end{enumerate}
For \texttt{194.1.1.1} the audit is more conservative: the new local library is internally validated, but still library-backed rather than independently benchmarked against an external BR/EBR source.

\section{Quotient Reinterpretation Formalism}
The crucial distinction is
\[
  BS/AI \;\; \text{as a raw algebraic quotient}
  \qquad \text{versus} \qquad
  X_{\mathrm{SI}} \;\; \text{as a finite indicator group}.
\]
If the Smith decomposition of the AI inclusion matrix inside a BS basis gives diagonal data
\[
  \operatorname{diag}(d_1,\dots,d_r),
\]
then the raw quotient has the structure
\[
  BS/AI \cong \mathbb{Z}^{(\operatorname{rank} BS-r)} \times \prod_{d_i>1} \mathbb{Z}_{d_i}.
\]
The free factor is not a finite indicator sector. Therefore any quotient with nonzero free rank must be interpreted carefully before using indicator language.

\section{10.4.1.31 Audit Result}
\subsection*{Single-group}
The BS layer is the least suspicious component. The with-planes nullity $8$ matches the stored BS basis size $8$, and the plane rank test is explicit. The AI layer also looks strong because all 15 families are enumerated and the local groups are only abelian or simple order-2 antiunitary case-a extensions. A fresh Smith decomposition of the stored AI basis coefficients reproduces the finite quotient
\[
  BS/AI \cong \mathrm{Z}_{2} \times \mathrm{Z}_{2}.
\]
Since there is no free part, the interpretation layer is comparatively stable here.

\subsection*{Double-group}
The raw quotient
\[
  BS/AI \cong \mathrm{Z}^{2} \times \mathrm{Z}_{2} \times \mathrm{Z}_{2} \times \mathrm{Z}_{2} \times \mathrm{Z}_{2}
\]
is mechanically well supported, and the corresponding report already notes that the quotient is not purely finite. The audit still ranks the interpretation layer as the weakest one, because the reporting artifacts continue to mix a free part with indicator language. The most honest reading is:
\[
  BS/AI = \underbrace{\mathbb{Z}^2}_{\text{free part}}
  \times
  \underbrace{\mathbb{Z}_2^4}_{\text{finite torsion part}}.
\]
Only the torsion sector is the natural finite-indicator candidate.

\section{194.1.1.1 Audit Result}
\subsection*{Single-group}
The BS layer is internally consistent but less settled than in \texttt{10.4.1.31} because it depends on synthetic boundary-point augmentation. The AI layer is substantially stronger after stage-2, but the completeness statement is still library-backed. The dominant problem is the interpretation of
\[
  BS/AI \cong \mathrm{Z}^{16}.
\]
Here the free rank is 16 and the finite torsion part is empty, so this is a raw free-dominated quotient, not a finite indicator group.

\subsection*{Double-group}
The same reinterpretation issue appears in the double branch:
\[
  BS/AI \cong \mathrm{Z}^{16}.
\]
The projective local library passes its internal validation, but the interpretation layer is still the first place to doubt. The new local library makes AI more credible than in stage-1, yet not so independently settled that it outranks the already obvious quotient-interpretation warning.

\section{Suspicion Ranking Table}
\begin{center}
\begin{tabular}{llll}
\toprule
case & top suspicion & second & least suspicion \\
\midrule
10.4.1.31 / single & AI / AI completeness & quotient interpretation & BS \\
10.4.1.31 / double & quotient interpretation & AI / AI completeness & BS \\
194.1.1.1 / single & quotient interpretation & AI / AI completeness & BS \\
194.1.1.1 / double & quotient interpretation & AI / AI completeness & BS \\
\bottomrule
\end{tabular}
\end{center}

\section{Implementation Mapping}
The machine-checkable outputs are:
\begin{itemize}
\item \texttt{two\_group\_correctness\_summary.json} for the four-case structured verdicts,
\item \texttt{audit\_10\_4\_1\_31\_correctness.md/json} for the per-group 10.4.1.31 audit,
\item \texttt{audit\_194\_1\_1\_1\_correctness.md/json} for the per-group 194.1.1.1 audit,
\item \texttt{reinterpretation\_194\_1\_1\_1\_raw\_vs\_finite.*} for the SG 194 raw-quotient reinterpretation,
\item \texttt{reinterpretation\_10\_4\_1\_31\_raw\_vs\_indicator.*} for the mixed-quotient reinterpretation on the reference group.
\end{itemize}

\section{Conclusion and Recommended Next Step}
The audit does \emph{not} find strong evidence that the BS layer is the main source of error in any of the four cases. The dominant warning sign is the quotient interpretation layer, especially whenever the reported quotient contains a free part. The cleanest corrective action is:
\begin{enumerate}
\item keep the raw quotients as algebraic outputs,
\item separate free and finite parts explicitly in the primary summaries,
\item reserve finite-indicator language for the torsion sector only,
\item treat the new SG 194 AI-completeness claims as library-backed until they are benchmarked independently.
\end{enumerate}

\end{document}
